% Generated draft: requires mathematical and source review.
\documentclass[11pt]{article}
\usepackage[margin=25mm]{geometry}
\usepackage{fontspec}
\setmainfont{Latin Modern Roman}
\usepackage{amsmath,amssymb,mathtools,unicode-math}
\setmathfont{Latin Modern Math}
\usepackage{enumitem,needspace,xurl,hyperref}
\setlength{\parindent}{0pt}
\setlength{\parskip}{5pt}
\setlength{\emergencystretch}{3em}
\newcommand{\N}{\mathbb{N}}
\newcommand{\Z}{\mathbb{Z}}
\newcommand{\Q}{\mathbb{Q}}
\newcommand{\R}{\mathbb{R}}
\newcommand{\C}{\mathbb{C}}
\newcommand{\F}{\mathbb{F}}
\newcommand{\A}{\mathbb{A}}
\newcommand{\PP}{\mathbb{P}}
\newcommand{\E}{\mathbb{E}}
\newcommand{\eps}{\varepsilon}
\newcommand{\dd}{\,\mathrm{d}}
\begin{document}
\section*{UnsolvedMath ID: 30003117}
\textbf{Source problem code:} catalogue entry
\subsection{Definitions and mathematical statement}
Let \(C\) be a complex algebraic curve, let \(\pi:A\to C\) be a family of abelian varieties of relative dimension at least two, and let \(s:C\to A\) be an algebraic section, so that \(\pi\circ s=\operatorname{id}_C\). For \(p\in C(\C)\), write \(A_p\) for the fiber and \(0_p\) for its identity. A point \(x\in A_p(\C)\) is torsion if \([n]_{A_p}x=0_p\) for some integer \(n\geq1\).

For a scheme-theoretic formulation, interpret the family as an abelian scheme over \(C\). Interpret an algebraic subgroup scheme as a closed subgroup scheme \(H\hookrightarrow A\) over \(C\), with the group operations inherited from \(A\); it need not be connected. Here ``proper subgroup scheme'' means \(H\neq A\), not a condition on the properness of its structure morphism. The supplied record does not explicitly specify these conventions or whether \(C\) is smooth, irreducible, or complete, so these aspects require confirmation against the source.

\emph{Question under this interpretation.} Suppose that
\[
 s(C)\not\subseteq H
 \qquad\text{for every closed subgroup scheme }H\subsetneq A\text{ over }C.
\]
Must the set
\[
 T_s=\left\{p\in C(\C):\exists n\in\Z_{\geq1},\ [n]_{A_p}s(p)=0_p\right\}
\]
be finite? The integer \(n\) may depend on \(p\). The question concerns arbitrary complex coefficients; the supplied formulation imposes no additional hypothesis of non-isotriviality or descent to \(\overline{\Q}\). This is a statement-only question, not an asserted finiteness theorem.
\subsection{Short English statement}
For a section of a family of abelian varieties of dimension at least two over a complex algebraic curve, ask whether avoiding every proper algebraic subgroup scheme forces torsion values to occur at only finitely many complex points of the base. The precise family and subgroup-scheme conventions require source review.
\subsection{Sources}
\begin{itemize}
\item Local UnsolvedMath registry, supplied record with permanent ID \(30003117\), title ``Torsion Specializations of Sections of Abelian Schemes,'' exact source problem code \texttt{catalogue entry}, in the Oberwolfach Reports Open Problems collection. No individual proposer is supplied; the proposed year is recorded as 2016.
\item Supplied source citation: \emph{Diophantische Approximationen} (2017), Oberwolfach Reports, pp.~1099--1187, DOI \texttt{10.4171/owr/2016/21}. \url{https://doi.org/10.4171/owr/2016/21}. The supplied background describes these reports as documenting mathematical workshops and their open-problem discussions.
\item Supplied literature source, ``OWR report'': \url{https://ems.press/content/serial-article-files/46625?nt=1}.
\item Supplied literature source, ``Corvaja--Masser--Zannier'': \url{https://doi.org/10.1007/s00208-017-1537-z}.
\item Supplied literature source, ``Masser--Zannier'': \url{https://doi.org/10.4310/ACTA.2020.v225.n2.a2}.
\item Literature assessment dated exactly \texttt{2026-08-21}: this is supplied dated curation, not independently verified. The registry labels the problem partially solved and reports that the curve and abelian-surface cases over complex numbers and several algebraic-number cases are proved, that later work treats arbitrary relative dimension over \(\overline{\Q}\), and that the full arbitrary-complex-coefficient statement remains unsettled. The supplied background calls this a dated literature triage derived from the supplied OWR curation and requests independent verification of time-sensitive status and cited claims.
\item The registry labels the statement \texttt{corrected\_verified}, explaining that mathematical content was preserved while notation, typography, or omitted definitions were repaired for self-contained publication. This supplied label does not constitute independent mathematical verification of this draft.
\end{itemize}
\end{document}
